% EDIT HERE: thesis metadata and accent colour. Then replace the sample chapters.
% One self-contained source; chapters can later be moved into \input files.
\documentclass[11pt,a4paper,oneside]{report}
\usepackage[T1]{fontenc}
\usepackage{tgpagella}
\usepackage[inner=30mm,outer=25mm,top=27mm,bottom=27mm]{geometry}
\usepackage{microtype,amsmath,amssymb,graphicx,booktabs,tabularx}
\usepackage{xcolor,titlesec,fancyhdr,enumitem}
\definecolor{accent}{HTML}{245C61}
\definecolor{muted}{HTML}{59676B}
\usepackage[colorlinks=true,allcolors=accent]{hyperref}
\newcommand{\thesistitle}{Thesis title}
\newcommand{\thesissubtitle}{A descriptive subtitle for your research}
\newcommand{\authorname}{Your name}
\newcommand{\university}{Your university}
\newcommand{\department}{Department / Faculty}
\newcommand{\degree}{Master's / Doctoral thesis}
\newcommand{\supervisor}{Supervisor name}
\hypersetup{pdftitle={\thesistitle},pdfauthor={\authorname}}
\linespread{1.12}
\setlength{\parindent}{0pt}\setlength{\parskip}{7pt}
\setlist[itemize]{leftmargin=*,itemsep=3pt,topsep=4pt}
\titleformat{\chapter}[display]{\bfseries\color{accent}}{\small\sffamily\MakeUppercase{\chaptertitlename}\ \thechapter}{8pt}{\Huge}
\titlespacing*{\chapter}{0pt}{0pt}{24pt}
\titleformat{\section}{\large\bfseries\color{accent}}{\thesection}{0.7em}{}
\pagestyle{fancy}\fancyhf{}
\fancyhead[L]{\small\color{muted}\nouppercase{\leftmark}}
\fancyfoot[L]{\small\color{muted}\authorname}
\fancyfoot[R]{\small\thepage}
\renewcommand{\headrulewidth}{0pt}\setlength{\headheight}{16pt}

\begin{document}
\hypersetup{pageanchor=false}
\begin{titlepage}
\raggedright
{\large\sffamily\color{accent}\university}\par
{\small\department}\par
\vspace{28mm}
{\small\sffamily\color{muted}\MakeUppercase{\degree}}\par
\vspace{8mm}
{\Huge\bfseries\thesistitle\par}
\vspace{6mm}
{\Large\color{muted}\thesissubtitle\par}
\vspace{12mm}{\color{accent}\hrule height 1.5pt}
\vspace{14mm}
{\large\authorname}\par
\vfill
\begin{tabular}{@{}ll@{}}
\textbf{Supervisor} & \supervisor\\[5pt]
\textbf{Programme} & Your programme\\[5pt]
\textbf{Submitted} & \today
\end{tabular}
\end{titlepage}
\pagenumbering{roman}
\hypersetup{pageanchor=true}
\chapter*{Abstract}
\addcontentsline{toc}{chapter}{Abstract}
Summarize the problem, the gap in existing knowledge and the main research
question. Describe your method, state the principal result and explain the
contribution. Write this section last so that it reflects the thesis you
actually produced.

This starter uses a study of research workflows as an illustrative thread.
Replace all example material, including the table and reference, with your
own evidence. Keep the abstract self-contained and follow your institution's
word limit and submission requirements.

\textbf{Keywords:} research question; methodology; evaluation; your field.
\vspace{12mm}
\section*{Acknowledgements}
Recognize supervision, collaboration, funding and personal support. Check whether
your institution requires a declaration of originality or other front matter.
\tableofcontents
\clearpage
\pagenumbering{arabic}
\chapter{Introduction and background}
\label{ch:introduction}
\section{Context and motivation}
Explain the problem in its broader setting. Identify who is affected, what is
currently possible and which difficulty motivates the research. Support factual
claims with appropriate sources, for example \cite{example}.

Move from the broad context to a specific gap in knowledge. Distinguish between
what earlier studies establish and what your thesis will investigate. Explain
why this gap matters academically or practically.
\section{Research question and objectives}
\textbf{Research question:} How does a revised workflow affect completion time
and accuracy compared with an existing process?
\begin{itemize}
\item Establish a defensible baseline and explicit evaluation criteria.
\item Develop or select an approach that addresses the identified gap.
\item Evaluate outcomes and discuss the limits of the evidence.
\end{itemize}
\section{Related work}
Organize the literature by themes or competing approaches. Compare their
assumptions, methods and findings rather than listing summaries. Conclude by
showing how the present study follows from that comparison.
\section{Scope and thesis structure}
State the boundaries of the study: population, setting, materials and outcomes.
Chapter~\ref{ch:methods} describes the method; Chapter~\ref{ch:findings} presents
and discusses the findings. Appendix~\ref{app:materials} contains supporting material.

\chapter{Research design and methods}
\label{ch:methods}
\section{Design and rationale}
Explain why your design can answer the research question. Identify the unit of
analysis, comparison conditions and controls. Where applicable, describe
recruitment, sampling and inclusion or exclusion criteria.
\section{Materials and procedure}
Document the materials, tools and versions used. Describe each stage in enough
detail for another researcher to reproduce it.
\begin{table}[htbp]
\centering
\caption{Example study protocol; adapt to your research.}
\begin{tabularx}{\linewidth}{@{}lX@{}}
\toprule
\textbf{Stage} & \textbf{Record and report}\\
\midrule
Preparation & Task specification, materials and selection criteria.\\
Collection & Conditions, observations and any protocol deviations.\\
Analysis & Metrics, uncertainty, exclusions and analysis code.\\
\bottomrule
\end{tabularx}
\end{table}
\section{Measures and analysis plan}
Define outcomes, units and the method of comparison before reporting results.
For completion time, one possible summary is
\begin{equation}
\Delta_t = \frac{t_{\mathrm{baseline}}-t_{\mathrm{revised}}}{t_{\mathrm{baseline}}}\times100\%.
\end{equation}
Explain your choice of statistical or qualitative analysis and the treatment of
missing observations. Distinguish confirmatory from exploratory analyses.
\section{Ethics and reproducibility}
Describe consent, privacy, data handling and any required approvals. State how
materials, data and code will be shared, or explain legitimate access restrictions.

\chapter{Findings, discussion and conclusions}
\label{ch:findings}
\section{Results}
Present the evidence before interpreting it. Report sample characteristics,
primary outcomes and uncertainty. Table~\ref{tab:outcomes} contains illustrative
values only; replace them with your measured results.
\begin{table}[htbp]
\centering
\caption{Illustrative comparison; these are not research findings.}
\label{tab:outcomes}
\begin{tabularx}{\linewidth}{@{}Xrr@{}}
\toprule
\textbf{Measure} & \textbf{Baseline} & \textbf{Revised}\\
\midrule
Completion time (minutes) & 24.0 & 18.0\\
Error rate (\%) & 8.0 & 5.0\\
\bottomrule
\end{tabularx}
\end{table}
\section{Interpretation and contribution}
Answer the research question directly. Explain how the evidence relates to the
objectives in Chapter~\ref{ch:introduction}, how it compares with earlier work
and what new understanding or capability follows from it.
\section{Limitations and future work}
Discuss threats to validity, sampling limitations and measurement assumptions.
Separate claims supported by your data from plausible explanations that need
further investigation. Propose concrete follow-up studies.
\section{Conclusion}
Summarize the answer, the contribution and its implications. This section should
stand on its own without introducing new evidence or claims.

\renewcommand{\bibname}{References}
\cleardoublepage\phantomsection
\addcontentsline{toc}{chapter}{References}
\begin{thebibliography}{9}
\bibitem{example} A. Author. \textit{Title of a relevant work}. Publisher, 2025.
% Replace this illustrative entry with real references for your research.
\end{thebibliography}
\appendix
\chapter{Supporting materials}
\label{app:materials}
\section{Reproduction checklist}
Provide instruments, extended results, additional proofs or implementation
details that support the main argument without interrupting it.
\begin{itemize}
\item Materials and instructions needed to repeat the study.
\item Data dictionary, units and anonymization decisions.
\item Software versions and analysis steps.
\item Additional tables or observations referenced in the main text.
\end{itemize}
\section{Example data record}
\begin{tabularx}{\linewidth}{@{}lX@{}}
\toprule
\textbf{Field} & \textbf{Meaning}\\
\midrule
Observation ID & Anonymous identifier for each recorded task.\\
Condition & Baseline or revised workflow.\\
Duration & Elapsed time, in minutes.\\
\bottomrule
\end{tabularx}
\end{document}
